\begin{center}
    {\Large\bfseries 7. MARCO REFERENCIAL}
\end{center}

\noindent\textbf{\large 7.1 Marco Histórico}\par\vspace{0.2cm}
La dependencia de los estudiantes universitarios hacia aplicaciones como Nequi, DaviPlata y Bancolombia es el resultado de una rápida transformación del ecosistema financiero colombiano que, en menos de una década, reemplazó el efectivo por los depósitos de bajo monto \cite{unisalle2026}. Este fenómeno histórico inició con la eliminación de las barreras bancarias tradicionales (papeleo, cuotas de manejo y sucursales físicas) a través de plataformas digitales, marcando un hito fundacional con el lanzamiento de Nequi en abril de 2016 y la expansión de DaviPlata \cite{unisalle2026}. Sin embargo, la adopción de estas tecnologías experimentó su aceleración definitiva con la pandemia de COVID-19, un evento que actuó como catalizador: el 61 \% de los colombianos disminuyó la frecuencia de uso del dinero físico y el volumen de pagos electrónicos aumentó un 70 \% ante las nuevas necesidades de bioseguridad \cite{cajamarca2021}.

Esta transición cualitativa provocó que, para el cierre de 2024, el 76,1 \% de la población adulta en Colombia ya utilizara depósitos de bajo monto, un producto que en poco tiempo logró que el 84,1 \% de sus cuentas se mantuvieran activas con un promedio de 92 operaciones anuales por usuario (casi el doble de la actividad registrada en las cuentas de ahorro tradicionales) \cite{unisalle2026}. La digitalización masiva de los micropagos mediante códigos QR y transferencias de número a número se afianzó aún más en la cotidianidad con la entrada en operación del sistema Bre-B en octubre de 2025 \cite{unisalle2026}. Es al final de esta evolución histórica donde la comunidad estudiantil de la UPTC, Sede Tunja, adaptó sus dinámicas de consumo —como el pago de fotocopias, alimentación en cafeterías y transporte urbano— a estas billeteras, llegando al punto de vulnerabilidad actual donde las caídas operativas de la infraestructura tecnológica paralizan su rutina académica y financiera \cite{restrepo2023habitos, martinez2023retos, pulido2024caracteristicas}.

\vspace{0.4cm}
\noindent\textbf{\large 7.2 Estado del Arte}\par\vspace{0.2cm}
\textit{Adopción y hábitos de uso de billeteras digitales en universitarios:} La literatura reciente coincide en la alta penetración de las plataformas financieras en la población joven. López y Vargas \cite{lopez2023adopcion} analizan la adopción desde el modelo UTAUT, destacando que los universitarios son altamente receptivos a la banca móvil por su facilidad de uso. En una línea similar, Restrepo y Gómez \cite{restrepo2023habitos} evalúan los hábitos de consumo en educación superior, coincidiendo en que las aplicaciones se utilizan predominantemente para transacciones inmediatas de bajo monto, como fotocopias y cafeterías. Ambos estudios concuerdan en la masificación de estas herramientas, pero se enfocan exclusivamente en los factores que motivan su adopción y uso en condiciones normales, sin contemplar escenarios de interrupción del servicio.

\textit{Infraestructura técnica e impacto operacional:} Desde una perspectiva más técnica, la literatura aborda las vulnerabilidades de estas plataformas. Martínez \cite{martinez2023retos} documenta los retos de disponibilidad en la infraestructura de la banca móvil colombiana, señalando la recurrencia de caídas en los servidores. Por su parte, Silva y Torres \cite{silva2022evaluacion} evalúan el impacto operacional de las aplicaciones de pago, concluyendo que las fallas de software reducen drásticamente la usabilidad percibida. Aunque ambos autores convergen en que la inestabilidad técnica es un problema latente, sus enfoques se centran en el rendimiento del software y en indicadores técnicos.

\textit{Identificación del vacío de investigación:} Al cruzar ambos grupos de estudios, se evidencia un vacío claro: la literatura actual documenta exhaustivamente la adopción tecnológica (desde el usuario) y los problemas de infraestructura (desde el software), pero ninguno de los estudios revisados aborda el cruce de ambas variables en un entorno microeconómico cautivo; es decir, no se evalúa el grado de parálisis o el impacto directo en las actividades cotidianas vitales (alimentación, movilidad, academia) de una comunidad estudiantil cuando ocurren estas caídas operativas. Ese es el vacío que aborda este proyecto, cuantificando empíricamente la dependencia funcional y la capacidad de contingencia en los estudiantes de la UPTC, Sede Tunja.

\vspace{0.4cm}
\noindent\textbf{\large 7.3 Marco Teórico}\par\vspace{0.2cm}
El marco teórico de esta investigación se sustenta en dos modelos conceptuales que permiten explicar tanto la adopción masiva de las plataformas financieras como el comportamiento de los estudiantes ante sus fallas operativas.

\begin{enumerate}
    \item \textit{Teoría Unificada de Aceptación y Uso de Tecnología (UTAUT):} Para abordar la primera y segunda sub-pregunta de investigación sobre el nivel de dependencia hacia Nequi, DaviPlata y Bancolombia, este estudio se apoya en el modelo UTAUT. Esta teoría explica por qué los usuarios adoptan un sistema de información basándose en la expectativa de rendimiento (la agilidad en los pagos), la expectativa de esfuerzo (facilidad de uso) y las condiciones facilitadoras (la infraestructura del campus). Bajo este modelo, la dependencia funcional que experimentan los estudiantes de la UPTC no es un fenómeno aleatorio, sino el resultado de que estas plataformas cumplen con alta eficiencia sus necesidades de transacciones cotidianas de bajo monto, generando un anclaje tecnológico profundo.

    \item \textit{Teoría de la Resiliencia de Sistemas y Dependencia del Medio:} Para abordar la tercera sub-pregunta, referente al impacto operacional y los mecanismos de contingencia frente a las caídas del servicio, se utiliza la intersección entre la Teoría de Sistemas (desde la resiliencia en la ingeniería) y la Teoría de la Dependencia del Medio. Esta perspectiva explica por qué una disrupción técnica en un nodo del sistema (el servidor del banco) genera un efecto en cascada que paraliza subsistemas físicos (alimentación, transporte urbano). Asimismo, permite interpretar, desde la elección racional, por qué los usuarios implementan (o fallan al implementar) mecanismos de contingencia, como el efectivo de reserva, buscando recuperar el equilibrio en sus actividades cotidianas ante la indisponibilidad de la infraestructura digital principal.
\end{enumerate}

\vspace{0.4cm}
\noindent\textbf{\large 7.4 Marco Conceptual}\par\vspace{0.2cm}
Para garantizar la precisión terminológica de esta investigación, se definen a continuación los conceptos centrales que estructuran el problema, organizados desde el fenómeno tecnológico general hasta las variables específicas de medición.

\begin{itemize}
    \item \textbf{Billetera digital y aplicación financiera móvil:} Herramienta de software instalada en un dispositivo móvil que permite a los usuarios almacenar dinero electrónicamente, realizar transferencias de fondos y ejecutar pagos digitales sin necesidad de intermediación física o dinero en efectivo \cite{bpmconsulting}. En el contexto de este estudio, este término se limita estrictamente a las plataformas Nequi, DaviPlata y la aplicación móvil de Bancolombia, excluyendo el análisis de pasarelas de pago web, criptomonedas o aplicaciones de entidades bancarias internacionales.
    
    \item \textbf{Dependencia funcional tecnológica:} Grado en el que un usuario necesita de manera indispensable un sistema de información para ejecutar tareas regulares, de modo que la ausencia del sistema imposibilita la acción \cite{saviasalud}. Para este proyecto, se define como la proporción de pagos diarios que el estudiante realiza exclusivamente a través de Nequi, DaviPlata o Bancolombia, y la incapacidad percibida para sustituir este medio de pago en su rutina universitaria inmediata.
    
    \item \textbf{Interrupción operativa del servicio (Caída o falla):} Evento en el cual la infraestructura tecnológica de la plataforma financiera (servidores, bases de datos o canales de comunicación) presenta latencias prolongadas, congelamiento de transacciones o inaccesibilidad total \cite{cyberhoot}. Para los fines de esta investigación, se consideran únicamente las fallas atribuibles a la entidad financiera, excluyendo interrupciones causadas por falta de datos móviles, descarga de batería o pérdida de contraseñas por parte del estudiante.
    
    \item \textbf{Actividades cotidianas (Contexto universitario):} Conjunto de transacciones de bajo monto y alta frecuencia necesarias para el desarrollo de la jornada presencial del estudiante. En este estudio, se delimitan específicamente a tres categorías dentro y en los alrededores de la UPTC Sede Tunja: 1) Alimentación (compra en cafeterías y restaurantes aledaños), 2) Apoyo académico (pago de fotocopias, impresiones y materiales), y 3) Movilidad (pago de transporte público urbano o colectivo). Se excluyen transacciones mayores como el pago de matrículas universitarias.
    
    \item \textbf{Mecanismo de contingencia:} Estrategia de respaldo o ruta alternativa que adopta un usuario u organización para mantener su operatividad frente a un evento de falla \cite{ibm_contingencia}. En este estudio, refiere a las acciones inmediatas que toma el estudiante ante la caída de su aplicación preferida, limitadas a: uso de efectivo de reserva guardado previamente, migración temporal a una billetera digital secundaria, préstamo de dinero por parte de un par, o la cancelación definitiva de la compra.
\end{itemize}